% Provenance (never printed): see sources/outline_sources.md (Year 2: tholeiitic, calc-alkaline
% and alkaline series; the TAS diagram; settings; isotopic signatures of mantle and crust).
% Classification from BGS RR 99-06 (after Streckeisen 1976, Le Bas et al. 1986, Le Maitre et
% al. 1989); series and data from Vermeesch and Pease (2021) and the GeoplotR package;
% curves and counts from figdata/bachelor-2/series.py; exercise and problem values in the
% Book 4 scratch calculators calc/ch09.py and calc/ch09pb.py.
\chapter{Series magmáticas y contextos tectónicos}\label{ch:b2:magmatic-series}

Dele a una petróloga un análisis químico y ningún mapa, y a menudo sabrá decir si la
lava salió en una dorsal, en una isla oceánica o sobre una placa que subduce. Las lavas
de Islandia y las de los volcanes de la cordillera de las Cascadas, en el oeste de Norteamérica, van ambas
del basalto a la riolita; analizadas por centenares, siguen dos trayectorias distintas al
evolucionar, una cada vez más rica en hierro, la otra no. Este capítulo empieza por
nombrar las rocas ígneas, a partir de sus minerales cuando se ven y de su
química cuando no; después las agrupa en series, familias de rocas
emparentadas por diferenciación, y vincula cada serie a un contexto tectónico, con los
diagramas que usan los petrólogos para leer un contexto en un análisis, y sus límites.
% ledger: vp21:series

\begin{recall}
\textbf{Del volumen de bachillerato}: el basalto, la andesita, la riolita, el gabro y el granito;
las rocas félsicas y máficas; las texturas volcánicas y plutónicas. \textbf{Del volumen de primer
año}: las dorsales oceánicas, las zonas de subducción y los arcos volcánicos, los puntos calientes y las grandes
provincias ígneas, los rifts continentales. \textbf{Del \cref{ch:b2:elements-in-the-earth}}:
los MORB y los OIB, los \omterm{def:b2:elements-in-the-earth:normalised}{diagramas normalizados} al \omterm{def:b2:elements-in-the-earth:bse}{manto primitivo} (de araña), los \omterm{def:b2:elements-in-the-earth:partition}{elementos compatibles} e
incompatibles. \textbf{Del \cref{ch:b2:isotope-geochemistry}}: el \omterm{def:b2:isotope-geochemistry:reservoirs}{manto
empobrecido} y los reservorios isotópicos enriquecidos. \textbf{Del
\cref{ch:b2:magmatic-differentiation}}: la \omterm{def:b2:magmatic-differentiation:crystallisation}{cristalización fraccionada}, la \omterm{def:b2:magmatic-differentiation:lld}{línea de
descenso del líquido}, los \omterm{def:b2:magmatic-differentiation:harker}{diagramas de Harker} y la ley de Rayleigh.
\end{recall}

\section{Nombrar las rocas ígneas}\label{sec:b2:magmatic-series:naming}

Una roca de grano grueso muestra sus minerales, y se nombra a partir de ellos. Una lava suele ser
de grano demasiado fino o demasiado vítrea para ello, y se nombra a partir de su análisis químico. Ambos
esquemas son convenciones internacionales: describen, y no implican nada sobre cómo
se formó la roca.
% ledger: bgs:tas-use

\begin{definition}[Composición modal; diagrama QAPF]\label{def:b2:magmatic-series:qapf}
La \emph{composición modal}\index{composición modal} de una roca es la proporción en
volumen de cada uno de sus minerales. El \emph{diagrama QAPF}\index{diagrama QAPF} nombra una
roca ígnea de grano grueso con menos de un \qty{90}{\percent} de minerales máficos a partir de
cuatro componentes modales recalculados a una suma de 100: cuarzo (Q), feldespato alcalino
(A), plagioclasa (P) y feldespatoides (F). Como el cuarzo y los feldespatoides no
cristalizan juntos, toda roca cae en uno de dos triángulos, QAP o FAP, unidos
por su lado común AP.
\end{definition}
% ledger: bgs:qapf

\begin{omfigure}
\centering
\begin{tikzpicture}[font=\scriptsize, x=1cm, y=1cm]
\def\L{6.4}
\pgfmathsetmacro{\H}{\L*0.8660254}
% upper point (q, p): q = Q fraction, p = P/(A+P)
\newcommand{\qapfup}[2]{({(#1)*\L/2 + (#2)*(1-(#1))*\L}, {(#1)*\H})}
\newcommand{\dn}[2]{({(#1)*\L/2 + (#2)*(1-(#1))*\L}, {-(#1)*\H})}
\draw[thick] (0,0) -- (\L/2,\H) -- (\L,0) -- (\L/2,-\H) -- cycle;
\draw (0,0) -- (\L,0);
\foreach \q in {0.9,0.6,0.2,0.05} \draw \qapfup{\q}{0} -- \qapfup{\q}{1};
\foreach \q in {0.1,0.6} \draw \dn{\q}{0} -- \dn{\q}{1};
\foreach \p in {0.1,0.65,0.9} \draw \qapfup{0.6}{\p} -- \qapfup{0}{\p};
\draw[dashed] \qapfup{0.6}{0.35} -- \qapfup{0.2}{0.35};
\draw \qapfup{0.2}{0.35} -- \qapfup{0}{0.35};
\foreach \p in {0.1,0.35,0.65,0.9} \draw \dn{0}{\p} -- \dn{0.1}{\p};
\foreach \p in {0.1,0.5,0.9} \draw \dn{0.1}{\p} -- \dn{0.6}{\p};
\node[anchor=south] at (\L/2,\H) {\normalsize Q};
\node[anchor=east] at (0,0) {\normalsize A};
\node[anchor=west] at (\L,0) {\normalsize P};
\node[anchor=north] at (\L/2,-\H) {\normalsize F};
\node at \qapfup{0.94}{0.5} {1a};
\node at \qapfup{0.75}{0.5} {1b};
\node[font=\tiny] at \qapfup{0.4}{0.05} {2};
\node at \qapfup{0.4}{0.225} {3a};
\node at \qapfup{0.4}{0.5} {3b};
\node at \qapfup{0.4}{0.775} {4};
\node[font=\tiny] at \qapfup{0.4}{0.95} {5};
\foreach \p/\n in {0.225/7*,0.5/8*,0.775/9*} \node at \qapfup{0.125}{\p} {\n};
\node[font=\tiny] at \qapfup{0.125}{0.05} {6*};
\node[font=\tiny] at \qapfup{0.125}{0.95} {10*};
\foreach \p/\n in {0.225/7,0.5/8,0.775/9} \node at \qapfup{0.025}{\p} {\n};
\foreach \p/\n in {0.225/7',0.5/8',0.775/9'} \node at \dn{0.05}{\p} {\n};
\node[font=\tiny] at \dn{0.3}{0.05} {11};
\node at \dn{0.35}{0.3} {12};
\node at \dn{0.35}{0.7} {13};
\node[font=\tiny] at \dn{0.3}{0.95} {14};
\node at \dn{0.75}{0.5} {15};
\node[anchor=east] at \qapfup{0.9}{0} {90};
\node[anchor=east] at \qapfup{0.6}{0} {60};
\node[anchor=east] at \qapfup{0.2}{0} {20};
\node[anchor=east] at \dn{0.1}{0} {10};
\node[anchor=east] at \dn{0.6}{0} {60};
\foreach \p/\n in {0.1/10,0.35/35,0.65/65,0.9/90} \node[anchor=south, fill=white, inner sep=0.5pt] at \qapfup{0.6}{\p} {\tiny\n};
\node[anchor=west, align=left, font=\scriptsize] at (\L+0.9,1.6) {Triángulo superior: \% Q\\en Q + A + P;\\
  rectas donde\\$100\,P/(A+P)$\\vale 10, 35, 65, 90};
\node[anchor=west, align=left, font=\scriptsize] at (\L+0.9,-1.7) {Triángulo inferior: \% F\\en F + A + P;\\
  rectas donde\\$100\,P/(A+P)$\\vale 10, 50, 90};
\end{tikzpicture}
\par\smallskip{\small La clasificación QAPF de las rocas ígneas de grano grueso. 1a
cuarzolita, 1b granitoide rico en cuarzo, 2 granito de feldespato alcalino, 3a sienogranito y
3b monzogranito (juntos, el granito), 4 granodiorita, 5 tonalita; 6*--10* las
variedades con cuarzo de 6--10: 6 sienita de feldespato alcalino, 7 sienita, 8 monzonita, 9 monzodiorita
o monzogabro, 10 diorita, gabro o anortosita; 6'--10' sus variedades con
foides; 11 sienita foídica, 12 monzosienita foídica, 13 monzodiorita o
monzogabro foídicos, 14 diorita o gabro foídicos, 15 foidolita. Los campos estrechos 6, 6', 10 y 10'
a lo largo de los bordes no están numerados. Según el esquema del British Geological Survey (basado
en Streckeisen 1976).}
\end{omfigure}
% ledger: bgs:qapf

El campo 10 contiene tres nombres. El \omterm{def:b2:magmatic-series:qapf}{diagrama QAPF} no separa la diorita, el gabro y la
anortosita, ni las rocas gabroicas entre sí; lo hacen otros diagramas, basados en la
plagioclasa y en los minerales máficos.

\begin{definition}[Diagrama de álcalis totales--sílice]\label{def:b2:magmatic-series:tas}
El \emph{diagrama de álcalis totales--sílice}\index{diagrama de álcalis totales--sílice} (diagrama
TAS) nombra una roca volcánica cuya moda no puede medirse a partir de sus porcentajes
en masa de \ce{Na2O + K2O} frente a \ce{SiO2}, una vez recalculado el análisis
al \qty{100}{\percent} sin agua ni dióxido de carbono. Está dividido
en quince campos por rectas cuyas intersecciones tienen coordenadas fijas
(la figura que sigue al próximo párrafo).
\end{definition}
% ledger: bgs:tas, bgs:tas-use

El diagrama tiene tres grupos de campos. A lo largo de la base, con pocos álcalis, están los
nombres de la serie común, del picrobasalto y el basalto a la riolita, separados en 45, 52,
57 y \qty{63}{\percent} de sílice y, para la riolita, por una recta inclinada. Encima, los campos S y la traquita contienen rocas más ricas en
álcalis con la misma sílice; los campos U, la fonolita y la foidita contienen rocas tan ricas en
álcalis, y tan pobres en sílice, que cristalizan feldespatoides. Algunos campos se
subdividen con un segundo criterio: un traquibasalto se llama hawaiita si $\ce{Na2O} - 2
\ge \ce{K2O}$, y traquibasalto potásico en caso contrario; una tefrita con más de un
\qty{10}{\percent} de olivino en su norma (véase más abajo) es una basanita; una riolita o una traquita cuya
razón molecular $(\ce{Na2O + K2O})/\ce{Al2O3}$ supera 1 se llama peralcalina. El esquema
exige análisis con menos de un \qty{2}{\percent} de agua y un \qty{0.5}{\percent} de
dióxido de carbono: la alteración añade agua y moviliza los álcalis, y un \omterm{def:b2:magmatic-differentiation:cumulate}{cumulado} no es un
líquido.
% ledger: bgs:tas, bgs:tas-use

\begin{omfigure}
\centering
\begin{tikzpicture}
\begin{axis}[omrecord, width=0.86\linewidth, height=8.6cm,
  xmin=37, xmax=80, ymin=0, ymax=16, xtick={37,41,...,77}, ytick={0,2,...,16},
  xlabel={\ce{SiO2} (\unit{\percent})}, ylabel={\ce{Na2O + K2O} (\unit{\percent})},
  legend style={font=\scriptsize, at={(0.98,0.98)}, anchor=north east},
  tick label style={font=\scriptsize}, label style={font=\small}]
\addplot[only marks, mark=*, mark size=0.55pt, blue!70!black, opacity=0.6] table[x=SiO2, y=alk] {figdata/out/bachelor-2-series-cath-th.dat};
\addlegendentry{Islandia}
\addplot[only marks, mark=*, mark size=0.55pt, red!75!black, opacity=0.6] table[x=SiO2, y=alk] {figdata/out/bachelor-2-series-cath-ca.dat};
\addlegendentry{Cascadas}
\foreach \f in {Pc,B,O1,O2,O3,R,S1,S2,S3,T,U1,U2,U3,Ph,F}
  \addplot[black, thin, forget plot] table[x=SiO2, y=alk] {figdata/out/bachelor-2-series-tas-\f.dat};
\node[font=\tiny] at (axis cs:43,1.6) {Pc};
\node[font=\tiny, fill=white, inner sep=0.5pt] at (axis cs:50.5,0.8) {B};
\node[font=\tiny] at (axis cs:54.5,2.2) {O1};
\node[font=\tiny] at (axis cs:60,2.2) {O2};
\node[font=\tiny] at (axis cs:66.5,2.2) {O3};
\node[font=\tiny] at (axis cs:74.5,11) {R};
\node[font=\tiny, fill=white, inner sep=0.5pt] at (axis cs:49,5.8) {S1};
\node[font=\tiny] at (axis cs:52.5,7.4) {S2};
\node[font=\tiny] at (axis cs:57.5,8.9) {S3};
\node[font=\tiny] at (axis cs:64.5,12) {T};
\node[font=\tiny] at (axis cs:44,6.6) {U1};
\node[font=\tiny] at (axis cs:48.7,9.4) {U2};
\node[font=\tiny] at (axis cs:52.5,11.8) {U3};
\node[font=\tiny] at (axis cs:57,14.6) {Ph};
\node[font=\tiny] at (axis cs:39,10) {F};
\end{axis}
\end{tikzpicture}
\par\smallskip{\small El diagrama TAS (según Le Bas y colaboradores, 1986, en la versión del
British Geological Survey) con 456 lavas y rocas intrusivas de Islandia (azul) y
580 de las Cascadas (rojo), óxidos tal como se analizaron. Pc picrobasalto, B basalto, O1 andesita
basáltica, O2 andesita, O3 dacita, R riolita; S1 traquibasalto, S2 traquiandesita
basáltica, S3 traquiandesita, T traquita o traquidacita; U1 tefrita o
basanita, U2 fonotefrita, U3 tefrifonolita, Ph fonolita, F foidita.}
\end{omfigure}
% ledger: bgs:tas, geoplotr:tas, data:b4-cath


\begin{definition}[Saturación en sílice]\label{def:b2:magmatic-series:saturation}
Un \omterm{def:b2:magma-emplacement:magma}{magma} o una roca está \emph{sobresaturado en sílice}\index{saturación en sílice} si tiene
más sílice de la que sus demás óxidos pueden combinar en minerales silicatados, de modo que
aparece cuarzo (u otro mineral de la sílice); \emph{subsaturado en sílice} si tiene demasiado
poca para convertir todos sus álcalis en feldespato, de modo que aparece en su lugar un feldespatoide (nefelina,
leucita), nunca junto con cuarzo; y \emph{saturado en sílice}
entre ambos casos.
\end{definition}

En una lava cuyos minerales no pueden contarse, la saturación se juzga sobre una \emph{norma}: un
cálculo estándar que reparte los óxidos del análisis entre minerales ideales
en un orden fijo, primero los álcalis en el feldespato, y da el exceso de sílice como
cuarzo o el déficit como feldespatoide. La más usada, la norma CIPW, separa los
dos tipos de basalto: un basalto con nefelina en su norma es un basalto alcalino; uno
sin ella es un basalto subalcalino, y entre estos los basaltos toleíticos llevan
hiperstena (ortopiroxeno) en la norma y, a menudo, piroxeno pobre en calcio en la moda.
El \cref{exo:b2:magmatic-series:11} aplica una versión simplificada.
% ledger: bgs:tas-use, bgs:series-terms

\section{Las series magmáticas}\label{sec:b2:magmatic-series:series}

Las lavas de un volcán, o de una provincia, rara vez son iguales, pero están emparentadas:
derivan de \omterm{def:b2:magma-emplacement:magma}{magmas} parentales semejantes por la diferenciación del
\cref{ch:b2:magmatic-differentiation}, y trazan juntas una línea de descenso.

\begin{definition}[Serie magmática]\label{def:b2:magmatic-series:series}
Una \emph{serie magmática}\index{serie magmática} es un conjunto de rocas ígneas de un mismo
ciclo magmático o de una misma provincia emparentadas por diferenciación a partir de \omterm{def:b2:magma-emplacement:magma}{magmas} parentales semejantes y
que comparten rasgos químicos y mineralógicos característicos. Se suelen
distinguir tres. La \emph{serie toleítica}\index{serie toleítica} parte de
basaltos saturados o ligeramente \omterm{def:b2:magmatic-series:saturation}{sobresaturados en sílice}, con hiperstena en la norma,
y muestra un marcado enriquecimiento en hierro al principio de su diferenciación. La
\emph{serie calcoalcalina}\index{serie calcoalcalina} va del basalto a la
andesita y después a la dacita y la riolita, con poco o ningún enriquecimiento en hierro, y con
minerales hidratados (hornblenda, biotita) en sus miembros evolucionados. La \emph{serie
alcalina}\index{serie alcalina} parte de basaltos con nefelina normativa y evoluciona
por los campos S y U del diagrama TAS hacia la traquita y la fonolita.
\end{definition}
% ledger: bgs:series-terms, bgs:tas-use

Los nombres describen conjuntos, no rocas aisladas: un basalto aislado puede pertenecer a cualquiera de
ellos, y los mismos nombres se han usado con significados distintos, así que una petróloga
debe decir qué entiende por ellos. La distinción entre la \omterm{def:b2:magmatic-series:series}{serie toleítica} y la
calcoalcalina se basa en el hierro, y el diagrama que mejor la muestra representa tres
grupos de óxidos.
% ledger: bgs:series-terms

\begin{definition}[Diagrama AFM]\label{def:b2:magmatic-series:afm}
Un \emph{diagrama AFM}\index{diagrama AFM} es un diagrama triangular cuyos vértices son
$A = \ce{Na2O + K2O}$, $F = \ce{FeO_T}$ (todo el hierro como \ce{FeO}) y $M = \ce{MgO}$, en
porcentaje en masa recalculado de modo que $A + F + M = 100$. Una línea trazada sobre él separa el
campo toleítico, del lado de F, del campo calcoalcalino.
\end{definition}
% ledger: vp21:series, geoplotr:afm

Un basalto que se diferencia pierde magnesio más deprisa que cualquier otra cosa, ya que el olivino y el
piroxeno son más ricos en magnesio que el líquido, y conserva sus álcalis, que son
incompatibles: su línea de descenso se aleja de M, hacia A. Lo que le ocurre al
hierro decide la serie. En un \omterm{def:b2:magma-emplacement:magma}{magma} toleítico el hierro solo sale con los
silicatos ferrosos ricos en magnesio, más despacio que el magnesio, de modo que la razón
$F/M$ aumenta y la tendencia sube primero hacia F antes de girar hacia A: la
\emph{tendencia de Fenner}. En un \omterm{def:b2:magma-emplacement:magma}{magma} calcoalcalino, más oxidado, el hierro también sale en la
magnetita, que contiene hierro férrico, y $F/M$ se mantiene casi constante mientras el líquido
va en línea recta hacia A: la \emph{tendencia de Bowen}. La figura siguiente muestra ambas con
rocas reales, las toleitas de Islandia y las lavas calcoalcalinas de las Cascadas. La
línea divisoria, trazada a ojo en 1971, las separa bien donde son basálticas; con
muchos álcalis las dos tendencias convergen y la línea dice poco.
% ledger: vp21:series, geoplotr:afm, data:b4-cath

\begin{omfigure}
\centering
\begin{tikzpicture}
\begin{axis}[name=afm, width=0.5\linewidth, height=0.5\linewidth, axis equal image,
  hide axis, xmin=-0.04, xmax=1.04, ymin=-0.06, ymax=0.92, clip=false]
\addplot[only marks, mark=*, mark size=0.5pt, blue!70!black, opacity=0.6] table[x=afm_x, y=afm_y] {figdata/out/bachelor-2-series-cath-th.dat};
\addplot[only marks, mark=*, mark size=0.5pt, red!75!black, opacity=0.6] table[x=afm_x, y=afm_y] {figdata/out/bachelor-2-series-cath-ca.dat};
\addplot[very thick, black] table[x=x, y=y] {figdata/out/bachelor-2-series-afm-line.dat};
\draw[thick] (axis cs:0,0) -- (axis cs:1,0) -- (axis cs:0.5,0.8660254) -- cycle;
\node[anchor=south] at (axis cs:0.5,0.866) {F};
\node[anchor=north east] at (axis cs:0,0) {A};
\node[anchor=north west] at (axis cs:1,0) {M};
\node[font=\scriptsize, fill=white, inner sep=1pt] at (axis cs:0.52,0.62) {toleítica};
\node[font=\scriptsize, fill=white, inner sep=1pt] at (axis cs:0.33,0.12) {calcoalcalina};
\end{axis}
\begin{axis}[omrecord, at={(afm.east)}, anchor=west, xshift=1.3cm,
  width=0.46\linewidth, height=0.46\linewidth,
  xmin=-2.5, xmax=2.2, ymin=-4.5, ymax=1.5,
  xlabel={$\ln(A/F)$}, ylabel={$\ln(M/F)$},
  legend style={font=\scriptsize, at={(0.98,0.98)}, anchor=north east},
  tick label style={font=\scriptsize}, label style={font=\small}]
\addplot[only marks, mark=*, mark size=0.5pt, blue!70!black, opacity=0.6] table[x=lnAF, y=lnMF] {figdata/out/bachelor-2-series-lr-th.dat};
\addlegendentry{Islandia}
\addplot[only marks, mark=*, mark size=0.5pt, red!75!black, opacity=0.6] table[x=lnAF, y=lnMF] {figdata/out/bachelor-2-series-lr-ca.dat};
\addlegendentry{Cascadas}
\addplot[thick, blue!50!black, forget plot] table[x=x, y=y_th] {figdata/out/bachelor-2-series-lr-fits.dat};
\addplot[thick, red!50!black, forget plot] table[x=x, y=y_ca] {figdata/out/bachelor-2-series-lr-fits.dat};
\end{axis}
\end{tikzpicture}
\par\smallskip{\small Izquierda: el \omterm{def:b2:magmatic-series:afm}{diagrama AFM} de las rocas de Islandia (azul) y de las Cascadas (rojo)
de la figura TAS, con la línea divisoria de Irvine y Baragar (1971). Derecha:
las mismas rocas en coordenadas de log-razones, donde cada serie se convierte en una banda casi
recta; las rectas son regresiones ortogonales, de pendientes $-1.10$ (Islandia) y $-0.72$
(Cascadas), que se cortan en $(-3.84, 2.12)$, como hallaron Vermeesch y Pease (2021).}
\end{omfigure}
% ledger: geoplotr:afm, data:b4-cath, vp21:model

De las rocas islandesas dentro del intervalo de la línea divisoria, 351 caen del lado
toleítico y 23 del otro; de las rocas de las Cascadas, 417 del lado calcoalcalino
y 57 del otro. Una curva es una frontera torpe, y un modelo sencillo explica
la forma de las tendencias. Supóngase que un \omterm{def:b2:magma-emplacement:magma}{magma} pierde sus álcalis, su hierro y su magnesio
a ritmos proporcionales a las cantidades que contiene, $\mathrm dX/\mathrm dt = -\lambda_X X$.

\begin{proposition}[Tendencias en log-razones]\label{prop:b2:magmatic-series:logratio}
Si $A$, $F$ y $M$ decrecen exponencialmente con constantes $\lambda_A$, $\lambda_F$ y
$\lambda_M$, el \omterm{def:b2:magma-emplacement:magma}{magma} se desplaza a lo largo de una recta en el plano de $\ln(A/F)$ y
$\ln(M/F)$, de pendiente $(\lambda_F - \lambda_M)/(\lambda_F - \lambda_A)$, mientras que en el triángulo AFM
se desplaza a lo largo de una curva.
\end{proposition}

\begin{proof}
Con $X = X_0 e^{-\lambda_X t}$, $\ln(A/F) = \ln(A_0/F_0) + (\lambda_F - \lambda_A)t$
y $\ln(M/F) = \ln(M_0/F_0) + (\lambda_F - \lambda_M)t$: ambas son lineales en $t$, luego
una es lineal en la otra, con la pendiente indicada. Las razones $A/F$ y $M/F$ fijan el
punto en el triángulo, cuyas coordenadas $A/(A + F + M)$, \dots\ no son lineales en $t$.
\end{proof}

La pendiente más fuerte de las rocas islandesas significa un $\lambda_F$ menor: el hierro sale
del \omterm{def:b2:magma-emplacement:magma}{magma} toleítico más despacio (\cref{exo:b2:magmatic-series:10}). Vermeesch y Pease mostraron que las dos
bandas de la figura anterior se cortan en un punto, que interpretaron como una hipotética fuente
común.
% ledger: vp21:model

La \omterm{def:b2:magmatic-series:series}{serie toleítica} y la calcoalcalina difieren también en su potasa. El
diagrama \ce{K2O}--\ce{SiO2} divide los basaltos, andesitas, dacitas y riolitas con
más de un \qty{48}{\percent} de sílice en tipos bajos en K, medios en K y altos en K mediante dos
rectas (figura siguiente). La mayor parte de ambos conjuntos son aquí medios en K: el
diagrama dice cuánta potasa lleva una serie, no qué serie es.
% ledger: bgs:k2o-sio2, data:b4-cath

\begin{omfigure}
\centering
\begin{tikzpicture}
\begin{axis}[omrecord, width=0.78\linewidth, height=6.2cm,
  xmin=45, xmax=79, ymin=0, ymax=5,
  xlabel={\ce{SiO2} (\unit{\percent})}, ylabel={\ce{K2O} (\unit{\percent})},
  legend style={font=\scriptsize, at={(0.02,0.98)}, anchor=north west},
  tick label style={font=\scriptsize}, label style={font=\small}]
\addplot[only marks, mark=*, mark size=0.55pt, blue!70!black, opacity=0.6] table[x=SiO2, y=K2O] {figdata/out/bachelor-2-series-cath-th.dat};
\addlegendentry{Islandia}
\addplot[only marks, mark=*, mark size=0.55pt, red!75!black, opacity=0.6] table[x=SiO2, y=K2O] {figdata/out/bachelor-2-series-cath-ca.dat};
\addlegendentry{Cascadas}
\addplot[thick, black, forget plot] table[x=x, y=y_low] {figdata/out/bachelor-2-series-k2o-lines.dat};
\addplot[thick, black, forget plot] table[x=x, y=y_high] {figdata/out/bachelor-2-series-k2o-lines.dat};
\draw[thick] (axis cs:48,0) -- (axis cs:48,0.3);
\node[font=\scriptsize, fill=white, inner sep=1pt] at (axis cs:60,3.7) {alto en K};
\node[font=\scriptsize, fill=white, inner sep=1pt] at (axis cs:76.5,2.2) {medio en K};
\node[font=\scriptsize, fill=white, inner sep=1pt] at (axis cs:60,0.35) {bajo en K};
\end{axis}
\end{tikzpicture}
\par\smallskip{\small Potasa frente a sílice para las rocas de Islandia y de las Cascadas, con
las rectas por (48, 0.3)--(68, 1.2) y (48, 1.2)--(68, 3.1) que separan los tipos bajos en K,
medios en K y altos en K (según Peccerillo y Taylor, 1976, tal como los adoptó el
esquema internacional), prolongadas aquí más allá de 68.}
\end{omfigure}
% ledger: bgs:k2o-sio2, data:b4-cath

¿Por qué difiere la oxidación? Los \omterm{def:b2:magma-emplacement:magma}{magmas} calcoalcalinos se forman sobre las zonas de subducción, donde
la placa que se deshidrata aporta especies oxidadas a la cuña del manto; los \omterm{def:b2:magma-emplacement:magma}{magmas}
cristalizan magnetita, y sus miembros evolucionados hornblenda o biotita. Los \omterm{def:b2:magma-emplacement:magma}{magmas}
toleíticos se forman en las dorsales, donde el manto se funde por descompresión
(\cref{ch:b2:mantle-melting}), y son más reducidos.
% ledger: vp21:series

\section{Series y contextos}\label{sec:b2:magmatic-series:settings}

Cada contexto de la tectónica de placas produce sus propios \omterm{def:b2:magma-emplacement:magma}{magmas}, porque cada uno funde el
manto a su manera: por descompresión bajo las dorsales y los puntos calientes, por aporte de agua
sobre las zonas de subducción (\cref{ch:b2:mantle-melting}).

Las \emph{dorsales oceánicas} emiten basaltos toleíticos, los MORB del
\cref{ch:b2:elements-in-the-earth}, pobres en \omterm{def:b2:elements-in-the-earth:partition}{elementos incompatibles} y con la
firma isotópica empobrecida de su fuente. Islandia, una dorsal sobre un punto caliente, añade
sus propias toleitas a esa lista, con una serie completa hasta la riolita.
% ledger: vp21:series

Las \emph{islas oceánicas} crecen por etapas. Un volcán hawaiano emite lavas alcalinas cuando es
joven, y después construye su escudo de basalto toleítico, que se estima en un 80 a un
\qty{95}{\percent} de su volumen final, en sus primeros 1 a 2 millones de años; al
alejarse del punto caliente sus lavas vuelven a ser alcalinas, y termina con una larga cola de
erupciones alcalinas escasas. Menos fusión da \omterm{def:b2:magma-emplacement:magma}{magmas} más alcalinos: los álcalis, al ser
incompatibles, se diluyen menos en una \omterm{def:b2:mantle-melting:fraction}{fracción de fundido} pequeña. Isotópicamente, los \omterm{def:b2:elements-in-the-earth:morb-oib}{basaltos de islas
oceánicas} muestrean los reservorios enriquecidos del \cref{ch:b2:isotope-geochemistry}.
% ledger: usgs14:hawaii-stages

Los \emph{arcos volcánicos} emiten \omterm{def:b2:magmatic-series:series}{series calcoalcalinas}. Sus \omterm{def:b2:elements-in-the-earth:trace}{elementos en trazas} llevan una
marca característica (figura siguiente): normalizados al \omterm{def:b2:elements-in-the-earth:bse}{manto primitivo}, los basaltos de
arco están enriquecidos en Rb, Ba, Th, U, Pb y Sr, elementos que la placa añade a la
cuña, y empobrecidos en niobio y tántalo respecto de sus vecinos. El basalto de arco
mediano del conjunto de datos tiene $\mathrm{Nb}/\mathrm{La} = 0.49$, frente a 0.97 para los MORB
y 1.23 para los \omterm{def:b2:elements-in-the-earth:morb-oib}{basaltos de isla oceánica}.
% ledger: data:b4-basalt-settings

\begin{omfigure}
\centering
\begin{tikzpicture}
\begin{axis}[omrecord, width=0.86\linewidth, height=6.4cm, ymode=log,
  xmin=0.5, xmax=22.5, ymin=1, ymax=300,
  xtick={1,...,22}, xticklabels={Rb,Ba,Th,U,Nb,Ta,La,Ce,Pb,Sr,Nd,Zr,Hf,Sm,Eu,Ti,Gd,Dy,Y,Er,Yb,Lu},
  ylabel={muestra / manto primitivo},
  legend style={font=\scriptsize, at={(0.98,0.98)}, anchor=north east},
  tick label style={font=\scriptsize}, label style={font=\small}]
\addplot[thick, mark=square*, mark size=1.4pt, orange!85!black] table[x=idx, y=oib] {figdata/out/bachelor-2-series-spider.dat};
\addlegendentry{basalto de isla oceánica}
\addplot[thick, mark=*, mark size=1.4pt, blue!70!black] table[x=idx, y=morb] {figdata/out/bachelor-2-series-spider.dat};
\addlegendentry{MORB}
\addplot[thick, mark=triangle*, mark size=1.7pt, red!75!black] table[x=idx, y=iab] {figdata/out/bachelor-2-series-spider.dat};
\addlegendentry{basalto de arco insular}
\end{axis}
\end{tikzpicture}
\par\smallskip{\small Composiciones medianas normalizadas al \omterm{def:b2:elements-in-the-earth:bse}{manto primitivo} de 241 MORB, 259
\omterm{def:b2:elements-in-the-earth:morb-oib}{basaltos de isla oceánica} y 256 basaltos de arcos insulares (la compilación de Vermeesch, 2006;
\omterm{def:b2:elements-in-the-earth:bse}{manto primitivo} de McDonough y Sun, 1995). Los basaltos de arco muestran el valle en Nb
y Ta y picos en Ba, Pb y Sr; los \omterm{def:b2:elements-in-the-earth:morb-oib}{basaltos de isla oceánica} están enriquecidos en todos los
elementos más incompatibles.}
\end{omfigure}
% ledger: data:b4-basalt-settings, ms95:pm-trace

Dentro de los continentes, el \omterm{def:b2:magma-emplacement:magma}{magma} basáltico encuentra una corteza gruesa que puede calentar y fundir,
y el resultado puede ser un vulcanismo con dos modas.

\begin{definition}[Vulcanismo bimodal]\label{def:b2:magmatic-series:bimodal}
El \emph{vulcanismo bimodal}\index{vulcanismo bimodal} es la erupción, en una misma provincia,
de dos grupos de \omterm{def:b2:magma-emplacement:magma}{magmas}, basálticos y riolíticos, con pocas composiciones intermedias
entre ellos.
\end{definition}

El vulcanismo a lo largo de la traza del punto caliente de Yellowstone es bimodal: en la llanura oriental
del río Snake, riolitas más antiguas con cerca de un \qty{75}{\percent} de sílice están cubiertas por
basaltos más jóvenes con cerca de un \qty{50}{\percent}. El hueco entre ellos es la clave: la
diferenciación de un basalto por sí sola habría dejado lavas intermedias, y un hueco
apunta a una segunda fuente para las riolitas, como la corteza que calienta el basalto.
% ledger: yvo23:bimodal

Las \emph{grandes provincias ígneas}, vistas en el volumen de primer año, emiten basalto en volúmenes
que ningún otro contexto iguala; su serie se lee con los mismos diagramas.

\section{Diagramas de discriminación}\label{sec:b2:magmatic-series:discrimination}

Las regularidades de la sección anterior invitan a leer un contexto directamente en un
análisis, y los petrólogos han trazado muchos diagramas para hacerlo.

\begin{definition}[Diagrama de discriminación]\label{def:b2:magmatic-series:discrimination}
Un \emph{diagrama de discriminación}\index{diagrama de discriminación} es un diagrama de
variables químicas, normalmente \omterm{def:b2:elements-in-the-earth:trace}{elementos en trazas} inmóviles, dividido en campos trazados
alrededor de las composiciones de rocas de contextos conocidos, que se usa para asignar un contexto a una
roca de origen desconocido.
\end{definition}

El clásico es el triángulo Ti--Zr--Y de Pearce y Cann (1973), trazado para basaltos.
Usa elementos que la alteración apenas moviliza, de modo que puede aplicarse a lavas antiguas
y metamorfizadas cuyo contexto se ha perdido. La figura siguiente lo pone a prueba con los 538 basaltos de la
compilación que tienen los tres elementos. La mayoría de los \omterm{def:b2:elements-in-the-earth:morb-oib}{basaltos de isla oceánica} caen en el
campo intraplaca D, 152 de 173, y la mayoría de los MORB en el campo B, 142 de 164; pero el campo B
es compartido con los basaltos de arco, y 101 de los 201 basaltos de arco caen también en él.
% ledger: geoplotr:tizry, data:b4-basalt-settings

\begin{omfigure}
\centering
\begin{tikzpicture}
\begin{axis}[width=0.62\linewidth, height=0.55\linewidth, axis equal image,
  hide axis, xmin=-0.04, xmax=1.04, ymin=-0.06, ymax=0.92, clip=false,
  legend style={font=\scriptsize, at={(1.0,0.95)}, anchor=north west}]
\addplot[only marks, mark=*, mark size=0.7pt, blue!70!black, opacity=0.7] table[x=x, y=y] {figdata/out/bachelor-2-series-tizry-morb.dat};
\addlegendentry{MORB}
\addplot[only marks, mark=square*, mark size=0.7pt, orange!85!black, opacity=0.7] table[x=x, y=y] {figdata/out/bachelor-2-series-tizry-oib.dat};
\addlegendentry{basalto de isla oceánica}
\addplot[only marks, mark=triangle*, mark size=0.9pt, red!75!black, opacity=0.7] table[x=x, y=y] {figdata/out/bachelor-2-series-tizry-iab.dat};
\addlegendentry{basalto de arco insular}
\foreach \f in {A,B,C,D}
  \addplot[thick, black, forget plot] table[x=x, y=y] {figdata/out/bachelor-2-series-tizry-\f.dat};
\draw[thick] (axis cs:0,0) -- (axis cs:1,0) -- (axis cs:0.5,0.8660254) -- cycle;
\node[anchor=south] at (axis cs:0.5,0.866) {Ti/100};
\node[anchor=north east] at (axis cs:0,0) {Zr};
\node[anchor=north west] at (axis cs:1,0) {Y $\times$ 3};
\draw[thin] (axis cs:0.33,0.36) -- (axis cs:0.22,0.46) node[anchor=south east, font=\small\bfseries, inner sep=1pt] {D};
\draw[thin] (axis cs:0.36,0.12) -- (axis cs:0.24,0.06) node[anchor=north east, font=\small\bfseries, inner sep=1pt] {C};
\draw[thin] (axis cs:0.56,0.37) -- (axis cs:0.76,0.50) node[anchor=west, font=\small\bfseries, inner sep=1pt] {A};
\draw[thin] (axis cs:0.53,0.21) -- (axis cs:0.72,0.13) node[anchor=north west, font=\small\bfseries, inner sep=1pt] {B};
\end{axis}
\end{tikzpicture}
\par\smallskip{\small El diagrama Ti--Zr--Y de Pearce y Cann (1973) con los basaltos
de contexto conocido. A toleitas de arco insular; B MORB, toleitas de arco insular y
basaltos calcoalcalinos; C basaltos calcoalcalinos; D basaltos intraplaca. Ti a partir de
\ce{TiO2}; los tres en ppm.}
\end{omfigure}
% ledger: geoplotr:tizry, data:b4-basalt-settings

\begin{method}[Uso de un diagrama de discriminación]\label{met:b2:magmatic-series:discrimination}
(1) Compruébese que la roca es del tipo para el que se construyó el diagrama, normalmente un basalto, y
no un \omterm{def:b2:magmatic-differentiation:cumulate}{cumulado} o una lava evolucionada, cuyos \omterm{def:b2:elements-in-the-earth:trace}{elementos en trazas} reflejan tanto la cristalización
como la fuente. (2) Para rocas alteradas o metamorfizadas, prefiéranse los elementos inmóviles durante la alteración (Ti, Zr, Y, Nb, las \omterm{def:b2:elements-in-the-earth:ree}{tierras raras}
pesadas). (3) Represéntese un conjunto de muestras, no una sola:
un único punto cerca de una frontera no decide nada. (4) Léase el resultado como una verosimilitud:
los campos se trazaron alrededor de la mayoría, no de todas, sus muestras de referencia, y algunos campos
son compartidos. (5) Confírmese con indicios independientes: otros diagramas, isótopos, el
contexto geológico.
\end{method}

El paso (4) puede cuantificarse (\cref{exo:b2:magmatic-series:8}). Si un basalto cae en el campo B, y los
tres contextos fueran igual de comunes, la probabilidad de que sea un MORB es solo de cerca de 0.62.
Los clasificadores estadísticos entrenados con tales compilaciones heredan el mismo límite: un basalto
guarda la memoria de su fuente y de su fusión, y dos contextos pueden producir \omterm{def:b2:magma-emplacement:magma}{magmas}
semejantes.
% ledger: data:b4-basalt-settings

\begin{inthelab}[Comprobar un análisis de roca total]
Antes de representar una roca en cualquiera de estos diagramas, los petrólogos comprueban el análisis. Los
óxidos mayores, medidos por ejemplo por fluorescencia de rayos X, deben sumar cerca de 100
con los volátiles: el agua y el dióxido de carbono, medidos directamente o juntos como la
pérdida por calcinación, la masa que pierde un polvo al calentarlo intensamente. Una gran pérdida por
calcinación señala una roca alterada, cuyos álcalis pueden haberse movilizado; el esquema del diagrama
TAS rechaza los análisis con más de un \qty{2}{\percent} de agua o un
\qty{0.5}{\percent} de dióxido de carbono. Los análisis se recalculan después a 100 sin
volátiles, el hierro se expresa como \ce{FeO} total o se reparte entre ferroso y férrico,
y los materiales de referencia analizados en la misma tanda se comparan con sus valores
certificados (\cref{ch:b2:field-lab-methods-2}).
\end{inthelab}
% ledger: bgs:tas-use

\section{Ejercicios}

\begin{exercise}[$\star$]\label{exo:b2:magmatic-series:1}
Nómbrense las siguientes rocas volcánicas en el diagrama TAS (\ce{SiO2}, \ce{Na2O}, \ce{K2O}
en \unit{\percent}, ya sin volátiles): (a) 49.60, 2.22, 0.51 (el basalto hawaiano
de referencia BHVO-2); (b) 58.5, 3.9, 1.6; (c) 72.0, 4.0, 4.2; (d) 47.0, 4.5, 2.5.
Para (d), nómbrense las dos rocas posibles y dígase qué decide entre ellas.
\end{exercise}
% ledger: usgs:bhvo2-major, bgs:tas

\begin{exercise}[$\star$]\label{exo:b2:magmatic-series:2}
Una lava da \ce{SiO2} 51.60, \ce{TiO2} 1.20, \ce{Al2O3} 16.10, \ce{FeO_T} 9.80,
\ce{MgO} 6.60, \ce{CaO} 8.20, \ce{Na2O} 3.50, \ce{K2O} 1.40, \ce{H2O} 1.40 y
\ce{CO2} \qty{0.20}{\percent}. (a) ¿Es aceptable el análisis para el diagrama TAS?
(b) Recalcúlese sin volátiles. (c) Nómbrese la roca antes y después del
recálculo. (d) ¿Es baja, media o alta en K?
\end{exercise}
% ledger: bgs:tas-use, bgs:k2o-sio2

\begin{exercise}[$\star$]\label{exo:b2:magmatic-series:3}
Nómbrese en el \omterm{def:b2:magmatic-series:qapf}{diagrama QAPF}: (a) una roca de grano grueso con un \qty{25}{\percent} de cuarzo,
un \qty{30}{\percent} de feldespato alcalino, un \qty{35}{\percent} de plagioclasa, un \qty{8}{\percent} de
biotita y un \qty{2}{\percent} de hornblenda; (b) otra con un \qty{5}{\percent} de cuarzo,
un \qty{5}{\percent} de feldespato alcalino, un \qty{55}{\percent} de plagioclasa y un
\qty{35}{\percent} de clinopiroxeno.
\end{exercise}
% ledger: bgs:qapf

\begin{exercise}[$\star$]\label{exo:b2:magmatic-series:4}
Clasifíquense como bajas, medias o altas en K, con las rectas de la figura \ce{K2O}--\ce{SiO2}:
(a) 50.0, 0.40; (b) 56.0, 1.50; (c) 62.0, 2.90; (d) 66.0, 1.00 (\ce{SiO2},
\ce{K2O}, \unit{\percent}).
\end{exercise}
% ledger: bgs:k2o-sio2

\begin{exercise}[$\star\star$]\label{exo:b2:magmatic-series:5}
(a) Calcúlense $A$, $F$ y $M$ para un basalto con \ce{FeO_T} 14.5, \ce{MgO} 4.5 y
\ce{Na2O + K2O} \qty{3.2}{\percent}, y para una andesita con 6.1, 3.8 y 5.05.
(b) La línea divisoria pasa por $(F, A, M) = (54.5, 23.0, 22.5)$, $(52.5, 30.0,
17.5)$ y $(49.6, 35.0, 15.4)$, el primero de los cuales es su punto más rico en hierro; el
lado toleítico es el de F. ¿De qué lado cae cada roca? (c) ¿Qué
mineral, al cristalizar pronto, mantiene la andesita donde está?
\end{exercise}
% ledger: geoplotr:afm, vp21:series

\begin{exercise}[$\star\star$]\label{exo:b2:magmatic-series:6}
El MORB mediano de la compilación tiene \qty{4}{ppm} de Nb y \qty{4.12}{ppm} de La, el
\omterm{def:b2:elements-in-the-earth:morb-oib}{basalto de isla oceánica} mediano 38 y 30.8, el basalto de arco mediano 3 y 6.1. Con los
valores del \omterm{def:b2:elements-in-the-earth:bse}{manto primitivo} Nb 0.658 y La \qty{0.648}{ppm}, calcúlense los valores normalizados
y $\mathrm{Nb_N}/\mathrm{La_N}$ de cada uno. ¿Cuál tiene una anomalía de niobio, y por qué?
\end{exercise}
% ledger: data:b4-basalt-settings, ms95:pm-trace

\begin{exercise}[$\star\star$]\label{exo:b2:magmatic-series:7}
El Ti es $0.5993$ de la masa del \ce{TiO2}. Calcúlense las coordenadas Ti--Zr--Y
(Ti/100, Zr, $3\mathrm Y$, recalculadas a 100) de tres basaltos: (a) \ce{TiO2}
\qty{1.6}{\percent}, Zr 110, Y \qty{32}{ppm}; (b) 2.8, 220, 29; (c) 0.8, 50, 20. Sitúense
en la figura y dígase qué permite concluir cada posición.
\end{exercise}
% ledger: geoplotr:tizry

\begin{exercise}[$\star\star$]\label{exo:b2:magmatic-series:8}
En la prueba Ti--Zr--Y de este capítulo, el campo B contiene 142 de 164 MORB, 4 de 173
\omterm{def:b2:elements-in-the-earth:morb-oib}{basaltos de isla oceánica} y 101 de 201 basaltos de arco. (a) ¿Qué fracción de cada grupo cae
en B? (b) Si los tres contextos fueran igual de comunes, ¿cuál es la probabilidad de que un
basalto que cae en B sea un MORB? (c) ¿Qué permite decir el campo B por sí solo?
\end{exercise}
% ledger: data:b4-basalt-settings

\begin{exercise}[$\star\star$]\label{exo:b2:magmatic-series:9}
Una riolita tiene \ce{Na2O} 5.5, \ce{K2O} 4.5 y \ce{Al2O3} \qty{10.5}{\percent} ($M$ =
61.98, 94.20 y \qty{101.96}{g/mol}). Calcúlese su razón molecular $(\ce{Na2O +
K2O})/\ce{Al2O3}$. (a) ¿Es peralcalina? (b) El feldespato alcalino contiene un
óxido alcalino por cada \ce{Al2O3}: ¿qué fracción de los álcalis no puede entrar en él?
\end{exercise}
% ledger: bgs:tas-use

\begin{exercise}[$\star\star\star$]\label{exo:b2:magmatic-series:10}
En el modelo de la proposición de las log-razones, tómense $\lambda_A = 0.4$ y $\lambda_M =
0.8$. (a) Calcúlese la pendiente de la tendencia para $\lambda_F = 0.5$ y para $\lambda_F =
0.7$. (b) Las rocas islandesas tienen una pendiente de $-1.10$ y las de las Cascadas
$-0.72$: ¿qué $\lambda_F$ implica cada una? (c) ¿Qué serie pierde su hierro más deprisa,
y qué mineral lo explica? (d) ¿Qué le ocurre a $F/M$ a lo largo de cada tendencia?
\end{exercise}
% ledger: vp21:model

\begin{exercise}[$\star\star\star$]\label{exo:b2:magmatic-series:11}
Una norma simplificada. Trabájese en moles de óxido. Póngase todo el \ce{K2O} en ortoclasa
(\ce{K2O.Al2O3.6SiO2}), todo el \ce{Na2O} en albita (\ce{Na2O.Al2O3.6SiO2}), el
\ce{Al2O3} restante con otro tanto de \ce{CaO} en anortita (\ce{CaO.Al2O3.2SiO2}), el
\ce{CaO} restante con otro tanto de $(\ce{Mg,Fe})\ce{O}$ en diópsido
(\ce{CaO.(Mg,Fe)O.2SiO2}) y el resto de $(\ce{Mg,Fe})\ce{O}$ en hiperstena
(\ce{(Mg,Fe)O.SiO2}). Si falta sílice, conviértase la hiperstena en olivino
(\ce{2(Mg,Fe)O.SiO2}); si aún falta, conviértase la albita en nefelina
(\ce{Na2O.Al2O3.2SiO2}). Aplíquese a (a) BHVO-2: \ce{SiO2} 49.60, \ce{Al2O3} 13.44,
\ce{FeO_T} 11.15, \ce{MgO} 7.257, \ce{CaO} 11.40, \ce{Na2O} 2.219, \ce{K2O}
\qty{0.513}{\percent}; (b) un basalto con 45.5, 15.0, 11.5, 8.5, 10.0, 3.8 y 1.5 (masas
molares \ce{SiO2} 60.08, \ce{Al2O3} 101.96, \ce{FeO} 71.84, \ce{MgO} 40.30, \ce{CaO}
56.08, \ce{Na2O} 61.98, \ce{K2O} \qty{94.20}{g/mol}). ¿Está cada uno saturado en sílice? ¿A
qué serie pertenece cada basalto?
\end{exercise}
% ledger: usgs:bhvo2-major, bgs:series-terms

\begin{exercise}[$\star\star\star$]\label{exo:b2:magmatic-series:12}
En la compilación, $\mathrm{Nb}/\mathrm{La} < 0.7$ se cumple en 116 de 141 basaltos de arco,
21 de 101 MORB y 3 de 137 \omterm{def:b2:elements-in-the-earth:morb-oib}{basaltos de isla oceánica} analizados para ambos. (a) Calcúlense las
tres fracciones. (b) Con proporciones iguales de los tres contextos, ¿cuál es la
probabilidad de que un basalto con $\mathrm{Nb}/\mathrm{La} < 0.7$ sea un basalto de arco? (c)
Rehágase para una región donde un \qty{60}{\percent} de los basaltos son MORB, un \qty{30}{\percent}
\omterm{def:b2:elements-in-the-earth:morb-oib}{basaltos de isla oceánica} y un \qty{10}{\percent} basaltos de arco. (d) ¿Qué lección encierra (c)
para el uso de los \omterm{def:b2:magmatic-series:discrimination}{diagramas de discriminación}?
\end{exercise}
% ledger: data:b4-basalt-settings

\section{Problema: un volcán, tres lavas}

\begin{problem}[{Problema de fin de semana --- de tres análisis al contexto de un volcán y
los cristales que dejó atrás}]
\label{pb:b2:magmatic-series:1}
Tres lavas de un estratovolcán (inventado), desde sus coladas de flanco más antiguas hasta su domo
más reciente, dan (\unit{\percent}):

\smallskip\noindent{\small
\begin{tabular}{@{}lrrrrrrrrrr@{}}
 & \ce{SiO2} & \ce{TiO2} & \ce{Al2O3} & \ce{FeO_T} & \ce{MgO} & \ce{CaO} & \ce{Na2O} & \ce{K2O} & \ce{H2O} & \ce{CO2}\\
lava 1 & 50.4 & 0.95 & 17.6 & 8.6 & 7.9 & 10.2 & 2.9 & 0.62 & 0.70 & 0.10\\
lava 2 & 58.6 & 0.80 & 17.4 & 6.1 & 3.8 & 6.9 & 3.9 & 1.15 & 1.10 & 0.05\\
lava 3 & 65.3 & 0.55 & 16.2 & 4.0 & 1.6 & 4.3 & 4.4 & 1.85 & 1.20 & 0.05\\
\end{tabular}}

\smallskip\noindent y en ppm: Rb 10, 20, 32; Th 1.1, 2.1, 3.4; Nb 2.0, 3.6, 5.6;
La 7, 13, 20; Sr 520, 480, 380; Zr 70, 125, 190; Y 18, 22, 24. \omterm{def:b2:elements-in-the-earth:bse}{Manto primitivo}: Nb
0.658, La 0.648, Th \qty{0.0795}{ppm}.

\noindent\textbf{Parte I --- Nombrar las lavas.}
\begin{enumerate}
\item Súmese cada análisis. ¿Son aceptables los tres para el diagrama TAS?
\item Recalcúlese la lava 1 sin volátiles: dense el factor, \ce{SiO2} y
\ce{Na2O + K2O}.
\item Hágase lo mismo con las lavas 2 y 3.
\item Nómbrense las tres lavas en el diagrama TAS.
\item ¿Cambiarían los nombres sin el recálculo?
\item Una roca de grano grueso con la composición de la lava 2 cristalizó en las raíces del
volcán. ¿Qué esquema la nombra, y en qué campo se esperaría encontrarla?
\end{enumerate}

\noindent\textbf{Parte II --- ¿Qué serie?}
\begin{enumerate}[resume]
\item Clasifíquese cada lava como baja, media o alta en K.
\item Calcúlense $A$, $F$, $M$ para cada una (sin volátiles).
\item La línea divisoria pasa por $(F, A, M) = (42.8, 17.1, 40.1)$, $(46.9, 18.1,
35.0)$, $(52.5, 30.0, 17.5)$, $(49.6, 35.0, 15.4)$, $(40.8, 50.0, 9.2)$ y $(37.9,
55.0, 7.1)$, con el campo toleítico del lado de F. Interpolando $F$ linealmente en
$A$ entre puntos vecinos, ¿de qué lado cae cada lava?
\item Calcúlese \ce{FeO_T}/\ce{MgO} para cada lava. ¿Hay enriquecimiento en hierro?
\item ¿A qué serie pertenece el volcán? ¿Qué mineral y qué condiciones
lo explican?
\end{enumerate}

\noindent\textbf{Parte III --- ¿Qué contexto?}
\begin{enumerate}[resume]
\item Normalícense Nb, La y Th de la lava 1 al \omterm{def:b2:elements-in-the-earth:bse}{manto primitivo}.
\item Calcúlese $\mathrm{Nb_N}/\mathrm{La_N}$ para la lava 1 y compárese con las medianas
del \cref{exo:b2:magmatic-series:6}.
\item Calcúlese $\mathrm{Th_N}/\mathrm{Nb_N}$. ¿Qué agente transportó el torio?
\item Calcúlense las coordenadas Ti--Zr--Y de la lava 1 (\ce{TiO2} del análisis sin
volátiles). ¿Qué dice su campo? ¿Por qué no representar las lavas 2 y 3?
\item Conclúyase: ¿dónde creció este volcán?
\end{enumerate}

\noindent\textbf{Parte IV --- ¿Cuánto cristalizó?}
\begin{enumerate}[resume]
\item La potasa no entra en ninguno de los primeros minerales. A partir del \ce{K2O} sin volátiles de
las lavas 1 y 2, calcúlese la fracción $F$ de líquido restante cuando se formó la lava 2, si
deriva de la lava 1 por \omterm{def:b2:magmatic-differentiation:crystallisation}{cristalización fraccionada}.
\item Hágase lo mismo con el rubidio y con el torio.
\item El circonio tiene un $D \approx 0.1$ global. Úsese la ley de Rayleigh para estimar $F$ a partir
de él.
\item El estroncio disminuye de la lava 2 a la lava 3. Calcúlese $F$ a partir de la potasa entre las lavas 2 y
3, y úsese para estimar el $D$ global del Sr en ese intervalo. ¿Qué mineral es
responsable?
\item De la lava 1 a la lava 3, ¿qué fracción del líquido cristalizó?
\item Si el \omterm{def:b2:magma-emplacement:magma}{magma} hubiera asimilado corteza granítica rica en potasa, ¿sería la estimación por la potasa
demasiado alta o demasiado baja?
\item Promédiense las cuatro estimaciones de la Parte IV para las lavas 1 a 2 y dese la masa de
cristales dejada atrás por cada \qty{100}{kg} de basalto para formar la andesita, el número clave
de este problema.
\end{enumerate}
\end{problem}
% ledger: bgs:tas-use, bgs:tas, bgs:k2o-sio2, geoplotr:afm, geoplotr:tizry, ms95:pm-trace
